\documentclass[12pt]{article}
\usepackage[utf8]{inputenc}
\usepackage[T1]{fontenc}
\usepackage{hyperref}
\title{Wagawaga Language Resources for K-12 Curriculum Development}
\author{Compiled by Claude Sonnet 5.5 (claude-sonnet-5-5) for Andr\'as Kornai}
\date{}
\begin{document}
\maketitle

\section{Introduction}
Wagawaga (ISO 639-3 wgb; also Baeaula) is an Austronesian (Oceanic, Suauic) language of the southeastern tip of Papua New Guinea, in Milne Bay Province. Sources give about 1,200 to 1,900 speakers. Some sources conflate it with Yaleba (Gamadoudou), which may be a separate language, and one Wikipedia page lists the code as ylb. Care is needed to avoid mixing the two. Very little material was found.

\section{Primary Educational Resources}
No primers, textbooks or school programmes were found in search results.

\section{Linguistic and Reference Materials}
\begin{itemize}
\item Wikipedia article ``Wagawaga language (New Guinea)''. \url{https://en.wikipedia.org/wiki/Wagawaga_language_(New_Guinea)}
\item Joshua Project language page for wgb. \url{https://joshuaproject.net/index.php/languages/wgb}
\end{itemize}

\section{Religious and Community Resources}
\begin{itemize}
\item Joshua Project profile of the Wagawaga (Deamuni) people. \url{https://joshuaproject.net/index.php/people_groups/15816}
\end{itemize}
No scripture was confirmed.

\section{Digital and Audio Resources}
None found.

\section{Curriculum Development Context}
The language is reported as stable and used by the whole community, but no teaching materials were found. A curriculum would require new documentation work, ideally with the community and Milne Bay education authorities.

\end{document}